\documentclass[11pt]{article}
\usepackage[T1]{fontenc}
\usepackage[utf8]{inputenc}
\usepackage{lmodern}
\usepackage[margin=1in]{geometry}
\usepackage{amsmath,amssymb,amsthm,mathtools}
\usepackage[hidelinks]{hyperref}
\hypersetup{
 pdftitle={Coarse endpoint concentration at every diverging scale for one-dimensional space-time random walks},
 pdfauthor={Alec Kriebel},
 pdfsubject={A self-contained concentration theorem with unresolved historical priority},
 pdfkeywords={space-time random environment, quenched concentration, coarse graining, coalescing coupling},
 pdflang={en}
}
\usepackage{microtype}
\newtheorem{theorem}{Theorem}
\newtheorem{lemma}[theorem]{Lemma}
\newtheorem{proposition}[theorem]{Proposition}
\newtheorem{remark}[theorem]{Remark}
\newcommand{\EE}{\mathbb E}
\newcommand{\PP}{\mathbb P}
\newcommand{\ZZ}{\mathbb Z}
\newcommand{\one}{\mathbf 1}
\newcommand{\cF}{\mathcal F}
\newcommand{\cP}{\mathcal P}
\newcommand{\norm}[1]{\left\lVert#1\right\rVert_1}
\DeclareMathOperator{\Var}{Var}
\title{Coarse endpoint concentration at every diverging scale\\
for one-dimensional space-time random walks}
\author{Alec Kriebel\\
\small Independent researcher\\
\small ORCID: \href{https://orcid.org/0009-0001-9320-500X}{0009-0001-9320-500X}}
\date{18 October 2026}
\begin{document}
\maketitle
\begin{abstract}
For a nearest-neighbor walk on $\ZZ$ in a fixed i.i.d.\ space-time
environment with transition probabilities strictly between zero and one,
we prove that its quenched and averaged endpoint laws agree, after
coarse graining into deterministic intervals of any diverging length,
outside an environment event of probability smaller than every inverse
power of time. No uniform ellipticity or quantitative rate of divergence
of the interval length is required. This is the affirmative
one-dimensional assertion contrary to Conjecture~1 in Berger's contribution
to Oberwolfach Report~38/2009. The proof combines replica overlap
estimates, environmental martingales, same-parity coalescence and
concentration over a short final segment.
\textbf{Priority qualification:} novelty and historical priority remain
unresolved. A bounded literature audit established no exact full
anticipation, but four relevant older original full texts could not be
accessed. We make no claim of firstness, independent discovery or that
the conjecture remained open until this work. Established late-layer
and coalescing methods, including those of Dunlap--Gu and
Barraquand--Rychnovsky, are explicitly credited.
\end{abstract}

\section{Statement, interpretation and prior work}

Let $\omega=(p_{s,x}:s\geq0,\ x\in\ZZ)$ have i.i.d.\ coordinates with
\begin{equation}\label{eq:ell}
             0<p_{s,x}<1\qquad\PP\text{-almost surely}.
\end{equation}
In a fixed environment the walk starts at $X_0=0$ and at time $s$ jumps
from $x$ to $x+1$ with probability $p_{s,x}$ and to $x-1$ otherwise.
Write $P^\omega$ for this quenched law,
$\mu_n^\omega(x)=P^\omega(X_n=x)$, and
$\alpha_n(x)=\EE\mu_n^\omega(x)$. The environment expectation and
probability are $\EE$ and $\PP$. We suppress $\omega$ in intermediate
calculations. For a deterministic partition $\cP$ of $\ZZ$ define
\[
 L_{n,\cP}=\sum_{I\in\cP}|\mu_n(I)-\alpha_n(I)|.
\]
This uses the full $\ell^1$ distance, twice the usual total variation
distance. For real $M>0$ and $b\in\mathbb R$, the intervals
$[b+jM,b+(j+1)M)$, $j\in\ZZ$, intersected with $\ZZ$, define a
partition; empty cells are harmless.

\begin{theorem}\label{thm:main}
For every fixed environment law satisfying \eqref{eq:ell}, every
deterministic sequence $M_N>0$ with $M_N\to\infty$, and every
$\varepsilon,K>0$, there are finite $C_{\varepsilon,K}$ and $N_0$
such that, for $N\geq N_0$ and every deterministic alignment $b_N$,
\begin{equation}\label{eq:main}
       \PP(L_{N,\cP_N}>\varepsilon)
                    \leq C_{\varepsilon,K}N^{-K},
 \qquad
 \cP_N=\{[b_N+jM_N,b_N+(j+1)M_N)\cap\ZZ:j\in\ZZ\}.
\end{equation}
The constants may depend on the fixed environment law.
The eventual threshold may also depend on the sequence $M_N$.
The bound is uniform in each deterministic alignment.
\end{theorem}

Here and below \emph{superpolynomially small} means $O(N^{-K})$ for
every fixed $K>0$, with its own constant and eventual threshold.
The theorem includes biased averaged walks, nearest-neighbor parity,
integer cells containing $M_N$ sites, real interval lengths and
arbitrarily slowly diverging nonmonotone scales.
It does not assert a simultaneous bound for environment-selected
partitions, a law-uniform bound as ellipticity degenerates, or a
result for static or correlated environments.

Berger's report \cite[pp.~2151--2152, Theorem~1 and Conjecture~1]{Berger}
states the analogous concentration assertion in transverse dimension
at least two and conjectures its failure in dimension one.
Theorem~\ref{thm:main} gives the affirmative assertion for its i.i.d.,
nearest-neighbor, elliptic one-dimensional model. A deterministic
elliptic law has $\mu_N=\alpha_N$; the theorem also covers every
nondegenerate law under \eqref{eq:ell}, so its conclusion is not based
on that trivial special case.

The difference-walk and quenched-mean mechanisms already occur in
Bal\'azs--Rassoul-Agha--Sepp\"al\"ainen \cite{BRAS}.
Independent particles and coupling in dynamical environments are
treated in Joseph--Rassoul-Agha--Sepp\"al\"ainen \cite{JRAS}.
Dunlap--Gu \cite{DG} isolate a short final layer in a
Gaussian-times-stationary-corrector local limit theorem for a continuum
stochastic flow with molecular Brownian noise.
Barraquand--Rychnovsky \cite[Section~1.4, Proposition~1.13 and
Section~4.5]{BR} prove a full-space beta-environment local limit
theorem and use final-layer decomposition and ordered coalescence.
Those mechanisms, as well as replica calculations, martingale exposure
and bounded differences, are established methods. All analytic inputs
needed here are proved below.

The modern distribution-valued fluctuation limits of
Drillick--Parekh \cite{DP} and Parekh \cite{Parekh} are also relevant.
The inspected statements concern different observables and convergence
topologies and do not state \eqref{eq:main}'s shrinking-cell
$\ell^1$ conclusion with a superpolynomial environment tail.
Nearest-neighbor parity is not a reason to exclude that framework:
Parekh \cite[Section~6.1, Remark~6.2]{Parekh} handles it by two-step
kernels on $2\ZZ$. A qualitative local or weak limit does not, by
itself, supply every inverse-power probability bound.

The bounded audit verified no earlier theorem with the full
quantifiers of \eqref{eq:main}; it did not establish novelty.
The complete original texts of Boldrighini--Minlos--Pellegrinotti
(1999) \cite{BMP1999}, Boldrighini--Pellegrinotti
(2001) \cite{BP2001}, Boldrighini--Minlos--Pellegrinotti
(2004) \cite{BMP2004}, and B\'erard (2004) \cite{Berard} could not be
accessed. Their metadata, abstracts or later primary discussions were
available, but do not exclude stronger estimates in the original
proofs. None of these four sources is represented as fully eliminated
as an anticipation. The accompanying priority ledger records exact
access and reading limits. Accordingly this note reports a verified
self-contained theorem with originality and historical priority
unresolved; it does not assert firstness, independent discovery,
a new method or present-day open status.

Put
\begin{equation}\label{eq:params}
 q=\EE p,\quad a=q(1-q),\quad
 \delta=\EE[p(1-p)],\quad \sigma^2=\Var(p)=a-\delta,\quad v=2q-1.
\end{equation}
Then $0<\delta\leq a\leq1/4$. The averaged walk has independent
sign increments of mean $v$ and variance $4a$. This follows by
successively averaging fresh time layers even when a path revisits
an old spatial location. Constants $C,c>0$ below may change from
line to line and depend on the fixed law. Let $\cF_s$ contain the
environment layers with time index less than $s$, and put
$I_s=\sum_x\mu_s(x)^2$.

\section{Replica returns and occupation moments}

\begin{lemma}\label{lem:overlap}
There is $C<\infty$ such that, for all deterministic integers $s,t\geq0$,
\begin{equation}\label{eq:condoverlap}
       \EE[I_{s+t}\mid\cF_s]\leq \frac{C}{\sqrt{t+1}}.
\end{equation}
For $L,k\geq1$ integers,
\begin{equation}\label{eq:moments}
 \EE\left[\left(\sum_{j=0}^{L-1}I_{s+j}\right)^k
                    \biggm|\cF_s\right]\leq k!C^kL^{k/2}.
\end{equation}
Consequently, for every fixed $\eta>0$, $L\leq N$ and deterministic
$s$, the event
$\sum_{j<L}I_{s+j}>N^\eta\sqrt L$ is superpolynomially unlikely.
\end{lemma}

\begin{proof}
Take two independent walks conditional on the same environment.
Their half-difference $D=(X-Y)/2$ is integral because the positions
have the same parity. On environmental averaging its transition
probabilities to each neighbor are $a$ away from zero and $\delta$
at zero; its holding probabilities are $1-2a$ and $1-2\delta$,
respectively. Indeed different current sites use independent fresh
variables, whereas at a shared site the probabilities of the two
opposite step pairs are both $\EE[p(1-p)]$. In particular,
\[
              \EE I_t=P_0(D_t=0)=:g_t .
\]

We derive a pointwise return bound. For $0<z<1$, the hitting
transform from $d\geq1$ to zero is $\rho(z)^d$, where the root in
$(0,1)$ solves
\[
 az\rho^2-[1-(1-2a)z]\rho+az=0.
\]
The other root exceeds one; the bounded solution of the hitting
recurrence is the stated one. First-return decomposition at zero gives
\begin{align}
 G(z):=\sum_{t\geq0}g_tz^t
 &=\frac{1}{1-(1-2\delta)z-2\delta z\rho(z)}\notag\\
 &=\frac{1}{
 (1-\delta/a)(1-z)+(\delta/a)
       \sqrt{(1-z)(1-(1-4a)z)}}\notag\\
 &\leq \frac{\sqrt a}{2\delta\sqrt{1-z}} .\label{eq:green}
\end{align}
The square root is positive on $(0,1)$; the last inequality uses
$\delta\leq a$ and $1-(1-4a)z\geq4a$.

To pass from this generating-function bound to individual coefficients,
use reversibility with weights $\pi(0)=a/\delta$ and $\pi(d)=1$
for $d\ne0$. The transition operator $P$ is a self-adjoint
contraction on $\ell^2(\pi)$: contraction follows from Jensen's
inequality and invariance of $\pi$. Every holding probability is
at least $1/2$, so $2P-\mathrm{Id}$ is also a self-adjoint Markov
contraction. Thus the spectrum of $P$ lies in $[0,1]$.
The spectral measure of $\one_{\{0\}}$, divided by $\pi(0)$,
is a probability measure $\nu$ with
$g_t=\int_{[0,1]}\lambda^t\,d\nu(\lambda)$.
Therefore $g_t$ decreases with $t$. For $n\geq1$ this implies
\[
 g_n\frac{1-z^{n+1}}{1-z}\leq G(z).
\]
Set $z=n/(n+1)$, note $z^{n+1}\leq e^{-1}$, and apply
\eqref{eq:green}. After including $n=0$, we obtain
$g_n\leq C/\sqrt{n+1}$.

We also need the backward return profile
$u_t(d)=P_d(D_t=0)$ to be maximal at zero.
It is symmetric. Write $b_t(d)=u_t(d)-u_t(d+1)$ for $d\geq0$.
The backward transition equation gives
\begin{align*}
 b_{t+1}(0)&=(1-2\delta-a)b_t(0)+a b_t(1),\\
 b_{t+1}(d)&=a b_t(d-1)+(1-2a)b_t(d)+a b_t(d+1),
                       &&d\geq1.
\end{align*}
All coefficients are nonnegative, since
$1-2\delta-a\geq1-3a\geq1/4$.
Initially $b_0(0)=1$ and all other $b_0(d)=0$.
Induction proves $u_t(d)\leq u_t(0)=g_t$.
Conditioning on $\cF_s$ now fixes the starting replica law
$\mu_s\otimes\mu_s$ and leaves later layers fresh, so
\[
 \EE[I_{s+t}\mid\cF_s]
    =\sum_{x,y}\mu_s(x)\mu_s(y)u_t((x-y)/2)\leq g_t.
\]
Every half-difference here is integral.

Let $A_L=\sum_{j<L}g_j\leq C\sqrt L$.
Expand the $k$th power in \eqref{eq:moments} and order its time
indices, overcounting by at most $k!$. In each resulting nonnegative
ordered sum, condition at its penultimate selected time.
The sum of conditional expectations of the last overlap is at most
$A_L$. Iterating removes all indices and costs at most $A_L^k$.
Repeated indices are included: at a zero time gap use $I_t\leq1=g_0$.
This proves \eqref{eq:moments} without independence of overlaps.
Markov's inequality gives $k!C^kN^{-k\eta}$ for the stated tail.
For each desired power choose a fixed sufficiently large $k$.
\end{proof}

\section{Environmental martingales and coarse expectation}

Let $T_tf(x)$ be the averaged expectation of $f$ for a $t$-step
walk started at $x$. For deterministic bounded $f$, the exposure
martingale for $\mu_n(f)$ is $\mu_s(T_{n-s}f)$, $0\leq s\leq n$.
Its next increment is exactly
\begin{equation}\label{eq:increment}
 D_s(f)=\sum_x\mu_s(x)(p_{s,x}-q)
       [T_{n-s-1}f(x+1)-T_{n-s-1}f(x-1)].
\end{equation}
If $h_s(x)$ denotes the bracket, independence within the next layer gives
\begin{equation}\label{eq:variation}
 |D_s(f)|\leq\sup_x|h_s(x)|,\qquad
 \EE[D_s(f)^2\mid\cF_s]
             =\sigma^2\sum_x\mu_s(x)^2h_s(x)^2.
\end{equation}
All sums at a fixed horizon have finite support.
For a $1/r$-Lipschitz $f$, the translation semigroup preserves its
Lipschitz constant; thus the increment bound is $2/r$ and the
predictable quadratic variation is at most
$Cr^{-2}\sum_{s<n}I_s$.

We use two elementary scalar concentration estimates, recorded here
to fix their conditional meaning. If martingale differences satisfy
$|D_j|\leq b$ and $Q=\sum_j\EE[D_j^2\mid\mathcal G_{j-1}]$, then
\begin{equation}\label{eq:freedman}
 \PP\left(\left|\sum_jD_j\right|\geq u,\ Q\leq w\right)
       \leq2\exp\left\{-\frac{u^2}{4(w+bu)}\right\}.
\end{equation}
Indeed $e^x\leq1+x+x^2$ for $|x|\leq1$ makes
$\exp(\lambda\sum_{i\leq j}D_i-\lambda^2Q_j)$ a nonnegative
supermartingale for $\lambda b\leq1$.
Markov's inequality on the joint event, with
$\lambda=u/[2(w+bu)]$, and the same argument for negative sums prove
\eqref{eq:freedman}. We never condition the martingale law on a
future small-variation event.

We will also use the following bounded-increment estimate.
For independent centered increments with range length at most two,
\[
       \EE e^{\lambda Z}\leq e^{\lambda^2/2}.
\]
To check this, the second derivative of the log moment-generating
function is a variance under an exponential tilt, bounded by one
for any variable supported on an interval of length two.
Integrate twice, using the zero mean.
The corresponding exponential supermartingale, stopped at the first
crossing and then bounded at the terminal time, gives
\begin{equation}\label{eq:maximal}
 \PP\left(\max_{t\leq h}|S_t|\geq u\right)
             \leq2e^{-u^2/(2h)}
\end{equation}
for centered partial sums of such independent increments.
The same argument, applied to conditional ranges of martingale
increments, gives the independent-input bounded-differences
inequality
\begin{equation}\label{eq:BD}
 \PP(F-\EE F>u)
       \leq\exp\left\{-\frac{2u^2}{\sum_jc_j^2}\right\}
\end{equation}
when altering input $j$ changes $F$ by at most $c_j$.
For completeness, exposing independent inputs yields martingale
increments of conditional range at most $c_j$; the preceding
log-moment argument after rescaling gives
$\EE(e^{\lambda D_j}\mid\mathcal G_{j-1})
 \leq e^{\lambda^2c_j^2/8}$.
Multiplication and optimization prove \eqref{eq:BD}.

\begin{lemma}\label{lem:expectation}
Uniformly in real interval lengths $M\geq1$, deterministic alignments,
$n\geq2$ and $B\geq1$,
\begin{equation}\label{eq:expectation}
 \EE L_{n,\cP}\leq
 C\sqrt{\left(\frac B M+n^{-1/2}\right)
       \left[1+\log(1+\min(M+1,\sqrt n))\right]}+\frac C{B^2}.
\end{equation}
In particular there is a deterministic bound $\rho(n,M)\leq2$,
uniform in alignment, with $\rho(n,M)\to0$ whenever both $n,M\to\infty$.
\end{lemma}

\begin{proof}
On its parity lattice the averaged kernel is the binomial mass
$b_t(j)=\PP(\operatorname{Bin}(t,q)=j)$.
For $\phi(\theta)=1-q+qe^{i\theta}$,
\[
 |\phi(\theta)|^2=1-4a\sin^2(\theta/2),\qquad
 |\phi(\theta)|^t\leq e^{-2at\theta^2/\pi^2},
                  \quad |\theta|\leq\pi.
\]
Fourier inversion, with and without the factor
$|e^{i\theta}-1|\leq|\theta|$, proves
\[
 \max_j b_t(j)\leq C(t+1)^{-1/2},\qquad
 \max_j|b_t(j+1)-b_t(j)|\leq C(t+1)^{-1}.
\]
Unimodality and zeros outside the support give
$\sum_j|b_t(j+1)-b_t(j)|=2\max_j b_t(j)$.
These are physical shifts by two, not one.

For $h_I(x)=T_t\one_I(x+1)-T_t\one_I(x-1)$, summing over the
contiguous parity sites of an interval telescopes to at most two
endpoint masses. Alternatively bound its at most $M+1$ site terms
by the adjacent-mass bound. The partition cannot increase the
$\ell^1$ norm of a signed measure, so
\begin{align*}
 \max_I|h_I(x)|
     &\leq C\min((t+1)^{-1/2},(M+1)(t+1)^{-1}),\\
 \sum_I|h_I(x)|&\leq C(t+1)^{-1/2},\\
 \sum_I h_I(x)^2
     &\leq C\min((t+1)^{-1},(M+1)(t+1)^{-3/2}).
\end{align*}
Orthogonality of \eqref{eq:increment} in time and
$\EE I_s=g_s$ therefore imply
\[
 \sum_I\Var(\mu_n(I))
  \leq C\sum_{s<n}(s+1)^{-1/2}
       \min((n-s)^{-1},(M+1)(n-s)^{-3/2}).
\]
For $s\leq n/2$ the sum is at most $Cn^{-1/2}$.
For $s>n/2$, bound $(s+1)^{-1/2}\leq Cn^{-1/2}$ and split the
remaining sum over $k=n-s$ at $k=(M+1)^2$.
This yields
\begin{equation}\label{eq:sumvar}
 \sum_I\Var(\mu_n(I))
       \leq\frac C{\sqrt n}
                 [1+\log(1+\min(M+1,\sqrt n))].
\end{equation}
The window $[vn-B\sqrt n,vn+B\sqrt n]$ intersects at most
$2B\sqrt n/M+2$ cells. Cauchy--Schwarz and \eqref{eq:sumvar}
bound the expected discrepancy on them by the square-root term
in \eqref{eq:expectation}. All other cells lie outside this window.
Their expected discrepancy is at most twice its averaged tail,
which is at most $C/B^2$ by the variance $4an$.
For fixed $B$, the square-root term tends to zero as $n,M\to\infty$:
if $M\leq\sqrt n$, use $\log(1+M)/M\to0$; otherwise use
$\log n/\sqrt n\to0$. Then let $B\to\infty$.
The infimum of the displayed bounds over $B\geq1$, capped at two,
defines a suitable $\rho$.
\end{proof}

\section{A polynomial auxiliary coarse scale}

The next estimate is used only at an auxiliary scale, not at the
possibly much slower target scale $M_N$.

\begin{proposition}\label{prop:coarse}
For fixed $\gamma>1/4$, $u>0$ and $K>0$,
\[
 \PP(L_{n,\cP}>u)=O(n^{-K})
\]
uniformly over integer cell lengths $M\geq n^\gamma$
and all deterministic alignments.
\end{proposition}

\begin{proof}
It suffices to treat $1/4<\gamma<1/2$, since a fixed exponent in
that range also covers all larger $\gamma$.
Put
\[
 d=\gamma-\tfrac14,\quad
 \beta=\tfrac14+\tfrac d2,\quad \eta=\tfrac d4,\quad
 r=n^\beta,\quad h=\left\lfloor\frac{n^{2\beta}}{(\log n)^2}\right\rfloor,
 \quad m=n-h.
\]
Eventually $1\leq h\leq n/2$. Average just the final $h$ layers:
$\nu_n=\EE[\mu_n\mid\cF_m]=\mu_mT_h$.
For any deterministic $|f|\leq1$, the gradient bound in the preceding
proof gives
\[
 |D_s(f)|\leq C/\sqrt h,\qquad
 Q_f\leq C\sum_{s<m}\frac{I_s}{n-s}.
\]
Divide these indices into deterministic consecutive dyadic blocks
$B_j$ on which
$2^jh\leq n-s-1<2^{j+1}h$.
Their lengths are at most $h_j=2^jh$, and there are $O(\log n)$
nonempty blocks. By Lemma~\ref{lem:overlap}, the common event
\[
 E_n=\left\{\sum_{s\in B_j}I_s\leq n^\eta\sqrt{h_j}
                         \text{ for all }j\right\}
\]
has superpolynomially small complement. A shorter last block can
be bounded by extending its deterministic interval to length $h_j$.
On $E_n$, $Q_f\leq Cn^\eta/\sqrt h$.
Equation~\eqref{eq:freedman}, applied to the joint event, yields,
for each fixed $t>0$,
\begin{equation}\label{eq:smoothtest}
 \PP(|\nu_n(f)-\alpha_n(f)|>t,\ E_n)
       \leq2\exp\{-c_t n^{\beta-\eta}/\log n\}.
\end{equation}

The window $W_n=[vn-\sqrt n\log n,vn+\sqrt n\log n]$ meets
$J\leq Cn^{1/2-\gamma}\log n+2$ cells. On those complete cells the
$\ell^1$ discrepancy is the maximum over their $2^J$ deterministic
sign tests, each bounded by one. Union bound \eqref{eq:smoothtest}
over the signs and charge $\PP(E_n^c)$ once.
Since
\[
       \beta-\eta-(1/2-\gamma)=5d/4>0,
\]
the negative exponential dominates $J\log2$.
The averaged mass outside $W_n$ is at most
$Ce^{-c(\log n)^2}$ by \eqref{eq:maximal} (or its terminal bound).
The expectation of $\nu_n(W_n^c)$ is this same mass.
Markov's inequality therefore controls the remaining cells with a
superpolynomial fixed-threshold tail. We have proved
\begin{equation}\label{eq:nusmooth}
       \PP(\norm{\nu_n(\cP)-\alpha_n(\cP)}>t)
                       \text{ is superpolynomially small}.
\end{equation}

We restore the final layers by controlling boundary crossings.
Represent the integer cells as real intervals with boundary grid
$b+M\ZZ$. Let
\[
 B_n=\{x\in\ZZ:\operatorname{dist}(x+vh,b+M\ZZ)\leq r\}.
\]
Choose the deterministic periodic majorant
\[
 f_B(x)=\max\{0,\min\{1,\,
                2-\operatorname{dist}(x+vh,b+M\ZZ)/r\}\}.
\]
It is $1/r$-Lipschitz, equals one on $B_n$, and vanishes outside
the radius-$2r$ boundary strips. Those strips are eventually
disjoint because $r/M\to0$ uniformly for $M\geq n^\gamma$.

For a unimodal probability sequence $(b_j)$ and integer $D\geq1$,
its mass on one residue class modulo $D$ is at most
$D^{-1}+2\max b_j$. To see this, discard at most two sampled
points next to the mode and compare each remaining sampled value
with the average of a disjoint $D$-term block toward the mode,
where the values are larger. The sum of these block averages
is at most $D^{-1}$.
A physical residue class modulo $M$ for $-m+2j$ corresponds,
when nonempty, to one binomial-index class modulo
$M/\gcd(2,M)$. Its mass is therefore at most
$2/M+2\max b_j$. The support of $f_B$ has at most $C(r+1)$
physical residues in one period, uniformly in the real shift.
The peak estimate gives
\[
       \EE\mu_m(f_B)\leq C(r+1)/M+C(r+1)/\sqrt m=o(1).
\]
For this Lipschitz test, \eqref{eq:variation} and
Lemma~\ref{lem:overlap} bound the predictable variation, off a
superpolynomially unlikely event, by
$Cn^{\eta+1/2-2\beta}$, with increment bound $2n^{-\beta}$.
For a fixed positive deviation, \eqref{eq:freedman} has exponent
at least
\[
        \frac{c}{n^{\eta+1/2-2\beta}+n^{-\beta}},
\]
which grows as a positive power since
$2\beta-1/2-\eta=3d/4>0$.
Thus $\mu_m(B_n)$ has a superpolynomial fixed-threshold tail.

Let
$D_n^\omega=P^\omega(|X_n-X_m-vh|>r)$.
By averaging fresh final layers and \eqref{eq:maximal},
$\EE D_n^\omega\leq Ce^{-cr^2/h}\leq Ce^{-c(\log n)^2}$.
Its fixed-threshold tail is superpolynomial by Markov's inequality.
Compare the actual endpoint and an averaged final endpoint with the
real position $X_m+vh$ under the same old law. If the start is outside
$B_n$ and the displacement error is at most $r$, the two compared
positions belong to the same cell. The coupling inequality gives
\[
 \norm{\mu_n(\cP)-\nu_n(\cP)}
       \leq4\mu_m(B_n)+2D_n^\omega+Ce^{-c(\log n)^2}.
\]
Half-open endpoint conventions cause no exception because equality
with a boundary is included in $B_n$.
Combining this bound with \eqref{eq:nusmooth} proves the proposition.
\end{proof}

\section{Final-segment smoothing and independent strips}

We now prove the all-scale assertion. Set, for large $N$,
\begin{equation}\label{eq:closureparams}
 \ell=\lfloor N^{3/4}\rfloor,\quad m=N-\ell,\quad
 H=\lceil N^{5/16}\rceil,\quad r=N^{3/8},\quad
 R=\lceil N^{3/8}\log N\rceil,\quad
 \eta=\kappa=1/32.
\end{equation}
Here $m\geq N/2$ eventually. All partitions and intervals introduced
below are deterministic.

\begin{lemma}\label{lem:oldmass}
For any prescribed family of at most $CN$ intervals, each containing
at most $10R+10$ integer sites,
\[
       \PP\left(\max_J\mu_m(J)>N^{-\kappa}\right)
                       \text{ is superpolynomially small}.
\]
\end{lemma}

\begin{proof}
For each interval choose a $[0,1]$ majorant equal to one on it,
tapering linearly to zero over distance $r$ on each side.
It is $1/r$-Lipschitz. The averaged binomial peak bound gives its
mean at most
$C(R+r+1)/\sqrt m=O(N^{-1/8}\log N)=o(N^{-1/32})$.
By Lemma~\ref{lem:overlap}, the single event
$\sum_{s<m}I_s\leq N^\eta\sqrt N$ has superpolynomially small
complement. On it the predictable variation for every majorant
is at most $CN^{-7/32}$ and the increment bound is $2N^{-3/8}$.
With $u=N^{-1/32}/2$, \eqref{eq:freedman} bounds each joint tail
by $2e^{-cN^{5/32}}$. Union bound over the prescribed intervals,
charging the common overlap failure only once.
\end{proof}

Write $K_\ell^\omega(x,\cdot)$ for the last-$\ell$-step kernel
in the fresh environment after time $m$.

\begin{lemma}\label{lem:coal}
For starts $x,y$ of the same parity,
\begin{equation}\label{eq:coalkernel}
 \EE_{\rm future}
     \norm{K_\ell^\omega(x,\cdot)-K_\ell^\omega(y,\cdot)}
            \leq C|x-y|/\sqrt{\ell+1}.
\end{equation}
Both sides may be capped by two.
\end{lemma}

\begin{proof}
In each fixed environment use independent transition coins until
the two paths meet and common coins thereafter. This preserves
each quenched marginal. Same-parity nearest-neighbor paths cannot
cross without meeting. Before meeting, environmental averaging
makes the positive half-gap a homogeneous lazy walk, killed at
zero, with neighbor probabilities $a,a$ and holding probability
$1-2a$. If $p_\ell(j)$ is its unrestricted mass from zero, the killed
kernel from $d\geq1$ to $z\geq1$ is
$p_\ell(z-d)-p_\ell(z+d)$. This follows either by reflection,
which leaves holding steps unchanged, or by checking its transition
recursion, zero boundary and initial values. Its survival probability is
\[
 \sum_{z\geq1}[p_\ell(z-d)-p_\ell(z+d)]
               =\sum_{j=1-d}^{d}p_\ell(j)
               \leq Cd/\sqrt{\ell+1}.
\]
The last bound follows by Fourier inversion:
$1-2a+2a\cos\theta$ is nonnegative for $a\leq1/4$ and at most
$e^{-c_a\theta^2}$ on $[-\pi,\pi]$.
The coupling inequality in full $\ell^1$ distance supplies factor
two, and $2d=|x-y|$. The case $x=y$ is zero.
\end{proof}

Condition on the old environment and write $\mu=\mu_m$.
In a fixed $H$-site partition $\cP_H$, replace the conditional
distribution within each cell by the conditional $\alpha_m$
distribution, keeping $\mu$'s cell mass:
\[
       \mu'(x)=\mu(J)\frac{\alpha_m(x)}{\alpha_m(J)},
                       \qquad x\in J,\ \alpha_m(J)>0.
\]
Null $\alpha_m$ cells also have zero $\mu_m$ mass: every reachable
parity site has positive $\alpha_m$ mass. Omit those cells.
Couple starts of $\mu$ and $\mu'$ within each cell; their distance
is less than $H$. Lemma~\ref{lem:coal} gives
\[
 \EE_{\rm future}\norm{\mu K_\ell-\mu'K_\ell}
                        \leq CH/\sqrt\ell,\qquad
 \norm{\mu'-\alpha_m}=L_{m,\cP_H}.
\]
The latter identity is an exact cellwise calculation.
Markov kernels and partitioning contract $\ell^1$ distance.
Moreover, convexity and Lemma~\ref{lem:expectation} give
\[
 \EE_{\rm future}
   \norm{\alpha_mK_\ell(\cP_N)-\alpha_N(\cP_N)}
                      \leq\rho(\ell,M_N).
\]
Indeed each start in the $\alpha_m$ mixture only translates the
deterministic target partition alignment. Thus, for
\[
       F=\norm{\mu_mK_\ell(\cP_N)-\alpha_N(\cP_N)},
\]
we have
\begin{equation}\label{eq:meantransfer}
 \EE[F\mid\cF_m]
     \leq L_{m,\cP_H}+CH/\sqrt\ell+\rho(\ell,M_N).
\end{equation}
Here $H/\sqrt\ell=O(N^{-1/16})$ and $\rho(\ell,M_N)\to0$
without a rate assumption on $M_N$.
Proposition~\ref{prop:coarse} applies at the fixed auxiliary
exponent $5/16$: $H\geq m^{5/16}$ and $m\geq N/2$.

It remains to concentrate $F$ over the future. From a start $x$
at time $m$, kill a path after its first step $t\in\{1,\ldots,\ell\}$
with
\[
                   |X_{m+t}-x-vt|>R,
\]
and let $K^\dagger(x,\cdot)$ be its surviving endpoint
subprobability kernel. Averaging and \eqref{eq:maximal} yield,
uniformly in $x$,
\begin{equation}\label{eq:killedmean}
 \EE_{\rm future}[1-K^\dagger(x,\ZZ)]
       \leq d_N:=Ce^{-cR^2/\ell}\leq Ce^{-c(\log N)^2}.
\end{equation}
This is a maximum over the entire segment.

Partition the fresh environmental coordinates into deterministic
moving strips
\[
 S_k=\{(t,y):0\leq t<\ell,\quad kR\leq y-vt<(k+1)R\},
                                  \qquad k\in\ZZ.
\]
Only finitely many reachable coordinates are needed in each strip.
These disjoint input vectors are independent and independent of
$\cF_m$, including when $v$ is irrational.
An alive path from $x$ can use only sites satisfying
$y-vt\in[x-R,x+R]$.
Consequently changing strip $k$ affects $K^\dagger(x,\cdot)$
only if
\[
            x\in J_k=[kR-R-1,(k+1)R+R+1]\cap\ZZ.
\]
This statement holds in both the original and the altered environment,
since the allowed tube is deterministic. No environmental variable at
a killed destination is consulted.

Put $w_k=\mu_m(J_k)$ and
$F^\dagger=\norm{\mu_mK^\dagger(\cP_N)-\alpha_N(\cP_N)}$.
Changing strip $k$ changes $F^\dagger$ by at most $2w_k$,
because each affected row is a subprobability measure and the norm
satisfies the reverse triangle inequality.
Each integer start belongs to at most six $J_k$: its admissible
real interval of indices has length $3+2/R\leq5$.
Therefore
\[
             \sum_kw_k\leq6,\qquad
             \sum_kw_k^2\leq6\max_kw_k.
\]
Only $O(N)$ strips can influence starts in $[-m,m]$;
choose a deterministic finite index range containing them all.
Their $J_k$ contain at most $10R+10$ integer sites.
Conditional on $\cF_m$, \eqref{eq:BD} now gives
\begin{equation}\label{eq:strip}
 \PP(F^\dagger-\EE[F^\dagger\mid\cF_m]>t\mid\cF_m)
        \leq\exp\{-t^2/(12\max_kw_k)\}.
\end{equation}
If the denominator is zero the function is constant and the tail
is zero. Lemma~\ref{lem:oldmass} bounds $\max_kw_k$ by
$N^{-1/32}$ off an old-environment event of superpolynomially
small probability. Thus each fixed positive upper deviation in
\eqref{eq:strip} has cost at most $e^{-c_tN^{1/32}}$.
The independent objects here are the strip vectors; no independence
of endpoint masses or overlapping starting cones is assumed.

Finally let $D=1-\mu_mK^\dagger(\ZZ)$.
Positivity of $K_\ell-K^\dagger$ gives $|F-F^\dagger|\leq D$.
By \eqref{eq:killedmean}, for every old environment,
\begin{equation}\label{eq:restore}
 \EE[D\mid\cF_m]\leq d_N,\quad
 \EE[F^\dagger\mid\cF_m]\leq\EE[F\mid\cF_m]+d_N,\quad
 \PP(D>t\mid\cF_m)\leq d_N/t .
\end{equation}

\begin{proof}[Proof of Theorem~\ref{thm:main}]
Fix $\varepsilon>0$.
For sufficiently large $N$,
$CH/\sqrt\ell+\rho(\ell,M_N)+d_N<\varepsilon/8$.
By Proposition~\ref{prop:coarse}, outside a superpolynomially
unlikely event, $L_{m,\cP_H}\leq\varepsilon/8$.
Equations~\eqref{eq:meantransfer} and \eqref{eq:restore} then give
$\EE[F^\dagger\mid\cF_m]\leq\varepsilon/4$.
Outside another superpolynomially unlikely event,
Lemma~\ref{lem:oldmass} supplies the strip weight bound.
If also $D\leq\varepsilon/4$, then $F>\varepsilon$ implies
$F^\dagger-\EE[F^\dagger\mid\cF_m]>\varepsilon/2$.
Equations~\eqref{eq:strip}--\eqref{eq:restore} therefore give,
for every fixed $K$,
\[
 \PP(F>\varepsilon)
    \leq O(N^{-K})+e^{-c_\varepsilon N^{1/32}}
                       +4d_N/\varepsilon.
\]
Both explicit errors are superpolynomial. Since $F=L_{N,\cP_N}$,
this proves \eqref{eq:main}. All uses of the target partition were
in norms, contraction and the real-interval expectation estimate;
the auxiliary integer partition imposes no restriction on $M_N$.
\end{proof}

\section{Scope, verification and disclosure}

The proof uses the fixed-law inequalities
$0<\delta\leq a\leq1/4$, fresh space-time independence and bounded
nearest-neighbor jumps. It uses neither inverse moments of $p$ or
$1-p$ nor a common positive lower support bound.
Same-parity coupling is sufficient because the old quenched and
averaged laws have the same reachable parity. An arrow law with
$p\in\{0,1\}$ has $\delta=0$ and is excluded; constants and
thresholds need not remain bounded as a family of elliptic laws
approaches such a law.

The asymptotic claim follows from the proofs above.
The accompanying verification package includes finite exact calculations
of replica transitions, overlap moments,
martingale identities, parity residue bounds, coalescing kernels
and strip influences;
such calculations check identities and edge cases and are not a
substitute for the analytic argument.

AI tools were used extensively in solving, drafting and checking
this work. Separate adversarial AI reviews of the underlying argument and exact
verification calculations informed this note.
It is unrefereed and has not undergone conventional human peer review.
The author is responsible for the statements and for any remaining
errors. The bounded priority audit leaves the explicit full-text
access gaps described in the introduction. The manuscript does not
certify a novel discovery or establish when the historical conjecture
was first resolved.

\begin{thebibliography}{99}
\bibitem{Berger}
N.~Berger,
\emph{Quenched Local Limit Theorem for Space-time Random Walk in
Random Environment}, pp.~2151--2152, in
\emph{Scaling Limits in Models of Statistical Mechanics},
Oberwolfach Report 38/2009.
\href{https://doi.org/10.4171/OWR/2009/38}
     {doi:10.4171/OWR/2009/38}.

\bibitem{BRAS}
M.~Bal\'azs, F.~Rassoul-Agha and T.~Sepp\"al\"ainen,
\emph{The random average process and random walk in a space-time
random environment in one dimension},
Commun.\ Math.\ Phys.\ \textbf{266} (2006), 499--545.
\href{https://www.math.utah.edu/~firas/Papers/rap.pdf}
     {Author manuscript}.

\bibitem{DG}
A.~Dunlap and Y.~Gu,
\emph{A quenched local limit theorem for stochastic flows},
J.\ Funct.\ Anal.\ (2022).
\href{https://doi.org/10.1016/j.jfa.2021.109372}
     {doi:10.1016/j.jfa.2021.109372}.

\bibitem{BR}
G.~Barraquand and M.~Rychnovsky,
\emph{Random walk on nonnegative integers in beta distributed
random environment},
Commun.\ Math.\ Phys.\ (2023).
\href{https://doi.org/10.1007/s00220-022-04536-1}
     {doi:10.1007/s00220-022-04536-1};
\href{https://arxiv.org/abs/2201.07270v3}
     {arXiv:2201.07270v3}, Sections~1.4 and~4.5.

\bibitem{JRAS}
M.~Joseph, F.~Rassoul-Agha and T.~Sepp\"al\"ainen,
\emph{Independent Particles in a Dynamical Random Environment},
in \emph{Probability and Analysis in Interacting Physical Systems},
Springer Proceedings in Mathematics \& Statistics \textbf{283}
(2019), 75--121.
\href{https://doi.org/10.1007/978-3-030-15338-0_4}
     {doi:10.1007/978-3-030-15338-0\_4}.
\bibitem{DP}
H.~Drillick and S.~Parekh,
\emph{Random walks in space-time random media in all spatial
dimensions: the full subcritical fluctuation regime} (2025).
\href{https://arxiv.org/abs/2510.22155v1}{arXiv:2510.22155v1}.

\bibitem{Parekh}
S.~Parekh,
\emph{Hierarchy of KPZ limits arising from directed random walk
models in random media},
Forum Math.\ Sigma (2026).
\href{https://doi.org/10.1017/fms.2026.10220}
     {doi:10.1017/fms.2026.10220}.

\bibitem{BMP1999}
C.~Boldrighini, R.~A.~Minlos and A.~Pellegrinotti,
\emph{Central limit theorem for a random walk in dynamical
environment: integral and local},
Theory Stoch.\ Process.\ \textbf{5(21)} (1999), nos.~3--4, 16--28.
Original full text unavailable in the present audit.

\bibitem{BP2001}
C.~Boldrighini and A.~Pellegrinotti,
\emph{$T^{-1/4}$-Noise for random walks in dynamic environment on $\ZZ$},
Mosc.\ Math.\ J.\ \textbf{1} (2001), no.~3, 365--380.
\href{https://doi.org/10.17323/1609-4514-2001-1-3-365-380}
     {doi:10.17323/1609-4514-2001-1-3-365-380}.
Original full text unavailable in the present audit.

\bibitem{BMP2004}
C.~Boldrighini, R.~A.~Minlos and A.~Pellegrinotti,
\emph{Random walks in quenched i.i.d.\ space-time random environment
are always a.s.\ diffusive} (2004).
\href{https://doi.org/10.1007/s00440-003-0331-x}
     {doi:10.1007/s00440-003-0331-x}.
Original full text unavailable in the present audit.

\bibitem{Berard}
J.~B\'erard,
\emph{The almost sure central limit theorem for one-dimensional
nearest-neighbour random walks in a space-time random environment}
(2004).
\href{https://doi.org/10.1239/jap/1077134669}
     {doi:10.1239/jap/1077134669}.
Original full text unavailable in the present audit.
\end{thebibliography}
\end{document}
